\section{Further Work}
\label{chap:further_work}
In this chapter, we will present our suggestions for what is next for Hover. The concept has a great potential for expansion. However, some improvements and further testing are left for the future because of the project’s time limitations.
Game Improvements
Our priority in this project has been to develop a working prototype according to our functional requirements (see Chapter \ref{chap:requirements}). 
Thus, features of lower priority were intentionally left out to focus on the more important requirements. Additionally, as a result of our experiment we became aware of new areas of improvements. The expansion possibilities in the application are categorized into four main areas: \textit{Improved social features}, \textit{Expansion of locations}, \textit{Further incentives for earning points}, and \textit{Technology Improvements}.


\textbf{Improved social feature}\\
Based on the feedback from the test participants (see Part \ref{part:results}), the most prominent area of improvement to the application is enhancing the social features. The current implementation does not include a social network. Thus, users cannot customize who they want to interact with. A social network would create a more personalized experience.

In addition to a customizable social network, additional features could enrich the social aspects of the game. An example is customizable clubs for friends or coworkers that can create better cohesion and relationships between players. These groups could then compete in various challenges, such as getting the most points in a certain period. Current implemented social features should also be extended to include this change. For instance, it should be possible to filter the leaderboard based on groups.

With an improved social network, the social experience could be enhanced by implementing more ways to interact. Comments on activities and direct messages are the most prominent examples of such improvements.


\textbf{Location expansion}\\
Hover zones are currently only found in the Trondheim area. Locations for other cities and places in Norway should be added to expand the game’s target group. The concept also has the potential to grow internationally by adding locations around the world. 

The location categories can also be extended by adding new ones that categorize locations not belonging to any of the four existing categories. The game could then be appealing to more players. An example of such a category is “outdoor life,” including parks, ski resorts, and hiking areas.


\textbf{Further incentives for earning points}\\
The results in Part \ref{part:results} also revealed that the incentives for spending points should be improved. One way to do this is to implement the possibility to buy in-game items. For example, the avatar could be extended with customization possibilities using items bought using points earned in the game. The more points the player spends, the cooler their avatar will look.

In the initial concept description in Chapter \ref{chap:new_ideas}, we presented avatar evolvability as an incentive for playing. By including this feature, we exploit Malone’s element of fantasy, where a player’s character depends on the number of points earned in a period. For instance, the more points a player gets in the exercise category, the more sporty the avatar will look. Moreover, if the player stops engaging in that category, their avatar will become less sporty.

Another solution is to add a possibility to use the points earned in the application on real-world rewards. Companies interested in people visiting their branches could be willing to provide prizes to the users. Examples include gift cards or various other sponsored items.

Additionally, new types of competitions and pre-defined challenges can further increase the incentives to earn points. An example of a new type of competition is “player of the week.” Pre-defined challenges are in-game challenges created by the application targeted at a specific group of users. They will lower the barrier for users to compete with friends.


\textbf{Technology improvements}\\
The last area of improvement is related to the technology and implementation details in the game. The process of tracking and publishing an activity (hovering) should be as simple as possible to incentivize usage. There are multiple ways the tracking process can be simplified.

The first suggestion is to improve the process of hovering together with a friend. Instead of using a shareable code, this feature can be implemented using Bluetooth. The application could then automatically boost points if players are within a certain distance of each other.

Another useful enhancement is to fine-tune the notifications and suggestions from the application. An auto-sensing feature starting a hover session when a user is inside a geofence could be implemented. Then if the player decides to track this activity, they will get points for the entire period inside the Hover zone.

In addition to this, the game should be tested on more device types. Due to the limited resources available in this project, the game has mainly been tested on personal devices. Unfortunately, both our devices run iOS, and thus Hover was primarily tested on this operating system during development. The application should be extensively tested on multiple phones from different manufacturers to ensure everything is working as expected on all devices.


\subsection{Additional User Testing of Hover}
The user testing phase of Hover has been influenced by multiple factors impacting the value of the research (as discussed in Chapter \ref{chap:retrospective}). The two most considerable ones were the current Covid-19 situation and the limited time available in this project. Consequently, the concept should be tested again when society is fully reopened, and people return to normal behavior. The duration of the testing should also be extended to potentially get a better sense of the participants’ change of habits. More extensive testing could produce other results than our testing and provide feedback on newly developed features.


\subsection{Summary}
This chapter has presented our suggestions for further work in the project. The topics discussed were related to how the application could be improved and suggesting a more extensive test phase under normal circumstances. The chapter is also the last one in the thesis, wrapping up the project by proposing ways of expanding and improving the concept.
